\documentclass[12pt, a4paper]{article}
\usepackage{../../../myconfig} % Gọi file cấu hình từ ngoài gốc vào

% ==========================================
% BẮT ĐẦU NỘI DUNG TÀI LIỆU
% ==========================================
\begin{document}

% Truyền 3 tham số: {Nhãn cam} {Chữ to chính giữa} {Chữ ở góc phải Header}
\titlebox{Chủ đề: Tích phân}{Tích phân}{Nguyên hàm - Tích phân: Tích phân}

% --- CÂU 1 ---
\questionLinesManual{0}{1}{Biết \( \displaystyle \int_{1}^{5} f(x) \, dx = 4 \). Giá trị của \( \displaystyle \int_{1}^{5} 3f(x) \, dx \) bằng}{
    \begin{tasks}(4)
        \task $7$
        \task $\dfrac{4}{3}$
        \task $64$
        \task $12$
    \end{tasks}
}

% --- CÂU 2 ---
\questionLinesManual{0}{2}{Biết \( F(x) = x^3 \) là một nguyên hàm của hàm số \( f(x) \) trên \( \mathbb{R} \). Giá trị của \( \displaystyle \int_{1}^{2} (2 + f(x)) \, dx \) bằng}{
    \begin{tasks}(4)
        \task $\dfrac{23}{4}$
        \task $7$
        \task $9$
        \task $\dfrac{15}{4}$
    \end{tasks}
}

% --- CÂU 3 ---
\questionLinesManual{0}{3}{Biết \( \displaystyle \int_{1}^{2} f(x) \, dx = 3 \) và \( \displaystyle \int_{1}^{2} g(x) \, dx = 2 \). Khi đó \( \displaystyle \int_{1}^{2} [f(x) - g(x)] \, dx \) bằng?}{
    \begin{tasks}(4)
        \task $6$
        \task $1$
        \task $5$
        \task $-1$
    \end{tasks}
}

% --- CÂU 4 ---
\questionLinesManual{0}{4}{Khẳng định nào trong các khẳng định sau đúng với mọi hàm \( f \), \( g \) liên tục trên \( K \) và \( a, b \) là các số bất kỳ thuộc \( K \)?}{
    \begin{tasks}(2)
        \task \( \displaystyle \int_a^b [f(x) + 2g(x)] \, dx = \int_a^b f(x) \, dx + 2\int_a^b g(x) \, dx \)
        \task \( \displaystyle \int_a^b \dfrac{f(x)}{g(x)} \, dx = \dfrac{\int_a^b f(x) \, dx}{\int_a^b g(x) \, dx} \)
        \task \( \displaystyle \int_a^b [f(x)g(x)] \, dx = \int_a^b f(x) \, dx \cdot \int_a^b g(x) \, dx \)
        \task \( \displaystyle \int_a^b f^2(x) \, dx = \left[ \int_a^b f(x) \, dx \right]^2 \)
    \end{tasks}
}

% --- CÂU 5 ---
\questionLinesManual{0}{5}{Cho \( \displaystyle \int_{-1}^{0} f(x) \, dx = 3 \), \( \displaystyle \int_{0}^{3} f(x) \, dx = 3 \). Tích phân \( \displaystyle \int_{-1}^{3} f(x) \, dx \) bằng}{
    \begin{tasks}(4)
        \task $6$
        \task $4$
        \task $2$
        \task $0$
    \end{tasks}
}

% --- CÂU 6 ---
\questionLinesManual{0}{6}{Giả sử \( I = \displaystyle \int_{0}^{\frac{\pi}{4}} \sin 3x \, dx = a + b \dfrac{\sqrt{2}}{2} \) (\( a, b \in \mathbb{Q} \)). Khi đó giá trị của \( a - b \) là}{
    \begin{tasks}(4)
        \task \( 0 \)
        \task \( -\dfrac{1}{6} \)
        \task \( -\dfrac{3}{10} \)
        \task \( \dfrac{1}{5} \)
    \end{tasks}
}

% --- CÂU 7 ---
\questionLinesManual{0}{7}{Tích phân \( \displaystyle \int_{0}^{1} (3x+1)(x+3) \, dx \) bằng}{
    \begin{tasks}(4)
        \task $12$
        \task $9$
        \task $5$
        \task $6$
    \end{tasks}
}

% --- CÂU 8 ---
\questionLinesManual{0}{8}{Tính tích phân \( I = \displaystyle \int_{1}^{e} \left( \dfrac{1}{x} - \dfrac{1}{x^2} \right) dx \).}{
    \begin{tasks}(4)
        \task \( I = \dfrac{1}{e} \)
        \task \( I = \dfrac{1}{e} + 1 \)
        \task \( I = 1 \)
        \task \( I = e \)
    \end{tasks}
}

% --- CÂU 9 ---
\questionLinesManual{0}{9}{Biết \( \displaystyle \int_{3}^{5} \dfrac{x^2 + x + 1}{x + 1} \, dx = a + \ln \dfrac{b}{2} \) với \( a, b \) là các số nguyên. Tính \( S = a - 2b \).}{
    \begin{tasks}(4)
        \task \( S = 2 \)
        \task \( S = -2 \)
        \task \( S = 5 \)
        \task \( S = 10 \)
    \end{tasks}
}
% --- CÂU 10 ---
\questionLinesManual{0}{10}{Cho \( \displaystyle \int_{0}^{1} \dfrac{1}{x^2 + 3x + 2} \, dx = a\ln 2 + b\ln 3 \), với \( a, b \) là các số hữu tỷ. Khi đó \( a + b \) bằng}{
    \begin{tasks}(4)
        \task $0$
        \task $2$
        \task $1$
        \task $-1$
    \end{tasks}
}



\end{document}
